\chapter{Conclusão}\label{cap:conclusao}

Esta dissertação partiu de uma observação simples sobre a reconstrução do estado de jogo. De entre os seis módulos que compõem a solução de referência do desafio SoccerNet, dois destacam-se pelo desempenho reduzido, a leitura do número de camisola e a atribuição da equipa e do lado do campo, e são precisamente os que menos atenção têm recebido na literatura. O trabalho concentrou-se nesses dois, mantendo todos os restantes exatamente como na \emph{baseline}. Esta delimitação não foi apenas uma questão de âmbito, foi também uma opção metodológica, uma vez que garante que qualquer variação de desempenho é atribuível ao que de facto se alterou, e não a uma reotimização global do sistema. É esta a diferença que separa este trabalho das soluções apresentadas ao desafio, que otimizam todos os módulos em simultâneo e não permitem, por isso, isolar o contributo de cada um.

Para a leitura do número, o reconhecimento ótico de caracteres da solução de referência deu lugar a um modelo multimodal de visão e linguagem, quantizado a 4 bits para caber nos 6~GB de memória gráfica disponíveis. A contribuição não está apenas na substituição do modelo, mas na forma como o problema lhe é apresentado. Em vez de lhe entregar o recorte completo do jogador, isola-se previamente a região dorsal por estimativa de pose, descartam-se os recortes sem informação utilizável e consolidam-se as leituras sucessivas de cada \emph{tracklet} por uma votação temporal em duas fases, na qual as respostas nulas pesam contra qualquer número candidato. Foram estes três passos, e não a escolha do modelo em si, que se revelaram determinantes.

Para a atribuição da equipa, corrigiram-se erros no agrupamento por aparência e no modo como cada grupo é associado a um lado do campo. O agrupamento passou a aferir a sua própria qualidade antes de ser aceite, recorrendo à cor da camisola quando as representações de reidentificação não separam as duas equipas. A atribuição do lado deixou de assentar num critério único e passou a resolver-se por uma cascata de quatro critérios de exigência decrescente, que termina num desempate pela posição do guarda-redes.

Os ganhos são claros. No conjunto de teste, a exatidão da leitura do número subiu de $18.6\%$ para $39.5\%$, com melhoria em $38$ dos $39$ vídeos comparados e sem uma única regressão, e a atribuição da equipa passou de $85.4\%$ para $97.8\%$, o que representa uma redução de cerca de $85\%$ nos \emph{tracklets} mal atribuídos. Na métrica oficial do desafio, o GS-HOTA subiu de $22.26\%$ para $44.26\%$, praticamente o dobro, obtido intervindo em apenas dois dos seis módulos.

A avaliação abrangeu a totalidade das $160$ sequências anotadas do conjunto de dados, e é neste ponto que reside a garantia mais importante do trabalho. O conjunto de treino nunca foi observado durante o desenvolvimento nem contribuiu para a afinação de qualquer limiar, funcionando por isso como conjunto de retenção. Os valores que aí se obtêm, $35.2\%$ na leitura do número e $42.57\%$ de GS-HOTA, ficam próximos dos do conjunto de teste, o que afasta a objeção mais previsível a um trabalho que afina limiares sobre parte dos dados de avaliação.

Delimitar o alcance destes ganhos ocupou uma parte substancial do trabalho. Três tentativas de melhorar a leitura do número, uma ponderação das leituras por qualidade, um realce de contraste e a apresentação conjunta de todas as vistas de um \emph{tracklet}, não produziram melhorias e em alguns casos degradaram resultados que já estavam corretos. As três chocaram com o mesmo obstáculo, o de que nos casos difíceis o recorte entregue ao modelo já não contém o número a ler. Para este módulo, a margem que resta encontra-se, por isso, nas etapas que o alimentam, e não na leitura nem na agregação.

No módulo de atribuição da equipa a situação é distinta, e convém enunciá-la sem a arredondar. O conjunto de validação fica $15$ pontos abaixo do de teste, e a análise conduzida na Secção~\ref{sec:desempenho-global} mostrou que essa diferença não pode, com as medições realizadas, ser imputada a uma causa em concreto. O GS-HOTA anula a similaridade de uma deteção sempre que um atributo não coincide, pelo que um valor baixo é compatível tanto com um detetor que não encontra os jogadores como com um erro sistemático de atributo sobre deteções corretas. A inspeção visual das sequências em causa revela a segunda dessas causas, com jogadores cujo equipamento se aproxima em tom do da equipa de arbitragem a serem classificados como tal. Fica por determinar o peso relativo de cada uma, e é este o trabalho imediato que o presente estudo deixa em aberto.

Na atribuição do lado do campo produziu-se ainda uma medição que a literatura não apresenta. Sobre a anotação de referência dos $160$ vídeos, verificou-se que o critério da posição média, adotado pela generalidade das soluções do desafio, se inverte em $9$ casos, ou seja em $5.6\%$ dos vídeos, ao passo que a posição do guarda-redes nunca contradiz a anotação nos $155$ vídeos em que é observável. O alcance deste resultado deve ser enunciado com precisão. O recurso ao guarda-redes como referência de lado não é original deste trabalho, pois já se encontra em soluções apresentadas ao desafio. O que aqui se acrescenta é a quantificação da frequência com que a suposição alternativa falha, que nenhuma dessas soluções apresenta, e a avaliação isolada do contributo de cada nível da cascata de decisão.

Essa medição sugere, aliás, uma limitação de conceção do próprio módulo aqui desenvolvido. O guarda-redes está anotado em $156$ dos $160$ vídeos, e nos nove vídeos do conjunto de treino em que a atribuição não acerta na totalidade dos \emph{tracklets} está presente em oito. O sinal existe, portanto, em quase todos os casos em que seria útil. Sucede que o desempate pelo guarda-redes ocupa o último lugar da cascata e só é acionado quando as posições médias dos dois grupos distam menos de $0.3$~m. Sempre que um critério anterior decide, ainda que incorretamente, o guarda-redes não chega a ser consultado. Num dos casos observados, as duas equipas encontram-se ambas na mesma metade do campo, separadas por menos de dez metros, ao passo que os respetivos guarda-redes distam mais de setenta. A informação que resolveria o caso estava disponível e não foi utilizada.

Quanto a trabalho futuro, identificam-se três linhas concretas, por ordem de prioridade.

A primeira é reavaliar o conjunto de validação com os atributos desativados na configuração do avaliador. Isso isola a componente puramente geométrica da deteção e permite repartir a queda observada entre a falha do detetor e o erro sistemático de atributo, questão que a Secção~\ref{sec:desempenho-global} deixa em aberto por não ser resolúvel com as medições disponíveis. É a experiência de menor custo e maior valor informativo de todas as que aqui se apontam.

A segunda é elevar o guarda-redes na cascata de decisão, associando-o a um dos grupos de jogadores em vez de o manter como desempate final. O sistema já calcula esta informação e descarta-a antes do agrupamento. Pelos números acima, a alteração teria efeito na maioria dos casos que ainda falham.

A terceira é alargar a avaliação a vídeos fora do conjunto de dados, de modo a caracterizar em que condições de captação cada módulo continua a funcionar. As sequências do SoccerNet-GSR provêm de uma única produção televisiva, o que limita o que se pode inferir sobre a generalização do sistema.

A estas acresce a direção que a análise das soluções mais competitivas aponta de forma independente, a de reconstruir a deteção e a calibração, precisamente os módulos que este trabalho manteve intactos. É aí que reside o salto para valores acima dos $60\%$, mas trata-se de um âmbito distinto do que aqui se delimitou.

Por último, vale a pena registar que a direção seguida coincide com a que o estado da arte veio a adotar. A formulação da identificação como tarefa multimodal, resolvida por um modelo de visão e linguagem, tornou-se prática corrente entre as soluções mais competitivas na edição de 2025, tendo sido aqui alcançada de forma independente e sob restrições de recursos consideravelmente mais severas.
